\documentclass[11pt,a4paper]{article}
\usepackage[T1]{fontenc}
\usepackage[utf8]{inputenc}
\usepackage{lmodern,microtype}
\usepackage[margin=25mm,headheight=15pt]{geometry}
\usepackage{amsmath,amssymb,amsthm,mathtools}
\usepackage{booktabs,longtable,array}
\usepackage{xcolor,enumitem,fancyhdr}
\usepackage{hyperref,xurl}
\hypersetup{colorlinks=true,linkcolor=blue!45!black,citecolor=blue!45!black,
 urlcolor=blue!45!black,pdftitle={Recovering signs and absorbing Pell quotients},
 pdfauthor={AI-assisted research prepared for Vladimir Reshetnikov}}
\setlength{\emergencystretch}{3em}
\setlist{itemsep=3pt,topsep=5pt}
\pagestyle{fancy}\fancyhf{}
\fancyhead[L]{\small Explicit Diophantine consistency equations}
\fancyhead[R]{\small 23 October 2026}\fancyfoot[C]{\thepage}
\newtheorem{theorem}{Theorem}[section]
\newtheorem{lemma}[theorem]{Lemma}
\newtheorem{proposition}[theorem]{Proposition}
\newtheorem{corollary}[theorem]{Corollary}
\theoremstyle{remark}\newtheorem{remark}[theorem]{Remark}
\newcommand{\ZZ}{\mathbb Z}
\newcommand{\NN}{\mathbb N}
\newcommand{\PP}{\mathbb Z_{>0}}
\newcommand{\Con}{\operatorname{Con}}
\newcommand{\ZFC}{\mathrm{ZFC}}
\newcommand{\Pos}{\operatorname{Pos}}
\newcommand{\sgn}{\operatorname{sgn}}
\newcommand{\refute}{\operatorname{Refute}_{\ZFC}}
\newcommand{\pin}{\texttt{beb0192cb8}}
% Frozen, verified primary metrics. Coefficients are fixed scalar leaves.
\newcommand{\PositiveOps}{87}
\newcommand{\PositiveExtra}{8}
\newcommand{\NineDegree}{1,573,266}
\newcommand{\NineOps}{504}

\title{Recovering signs and absorbing Pell quotients\\[5pt]
\large New explicit Diophantine equations for syntactic ZFC consistency}
\author{Prepared for Vladimir Reshetnikov\\
\small AI-assisted mathematical research for the ProveIt collection}
\date{23 October 2026}
\begin{document}
\maketitle
\begin{abstract}
We construct an explicit quartic polynomial in \textbf{42 unrestricted integer
unknowns}, improving the 45-integer quartic in the inspected ProveIt snapshot.
It is the sum of the squares of 27 quadratic residuals. The new mechanism
recovers the signs of two coupled Pell coordinates from their common sign:
one ternary positivity certificate replaces two, while the harmless canonical
restriction $w=4u^2$ supplies a stronger modulus that excludes the negative
branch. No extra lifting variable is introduced.

A second construction gives an equation in \textbf{14 positive unknowns},
of \textbf{exact degree 34}, with a \textbf{\PositiveOps-operation} evaluator.
It improves the prior degree-52 point at that variable count. An existing
square equation carries a divisibility condition, and a positive main-root
coordinate absorbs the index quotient. A first-index lower bound obtained
before reconstructing that quotient is essential to the proof.

The additional nine-positive construction in Section~\ref{sec:deflated-nine} gives
another degree/circuit-size tradeoff. All equations are ordinary polynomials:
their exponents are fixed integers and their coefficient data are explicitly
bound. At the ordinary mathematical proof level of the named inherited
compiler and calibration theorems, ZFC proves that each equation has no
solution on its stated domain exactly when ZFC is syntactically consistent.
The new reductions are proved here; finite programs check their literal
arithmetic data. No new end-to-end proof-assistant certificate or global
minimality theorem is asserted.
\end{abstract}
\clearpage\tableofcontents\clearpage

\section{Results and the meaning of simplicity}
\label{sec:results}

An effective consistency equation must contain a fully specified integer
polynomial, a declared domain for every unknown, and a proof of both
directions of its refutation/solution equivalence. Writing a universal
polynomial with an unspecified program index does not meet that standard.
Neither does using a coefficient chosen according to the truth of
$\Con(\ZFC)$. The constructions below use only finite, reproducible
coefficient recipes; none searches for a refutation while defining a
coefficient.

The comparison snapshot is
\begin{center}\small
\texttt{beb0192cb8dc36f78c369c080f2adfbc1887a486}.
\end{center}
This commit includes sources 54--59, the newly placed eighth wave. The
collection's top-level comparison still describes the preceding wave, so
we use the individual new packages as the actual comparison points
\cite{collection,p55,p57,p58,p59}.

\begin{table}[htbp]\centering\small
\begin{tabular}{@{}llrrr@{}}
\toprule
Construction & Domain & Unknowns & Exact degree & Binary operations\\
\midrule
Source 57 quartic & $\ZZ$ &45&4&not optimized\\
New $Q_{42}$ & $\ZZ$ &\textbf{42}&\textbf{4}&not optimized\\
Sources 55/59 & $\PP$ &14&52&79\\
New $P_{14}$ & $\PP$ &\textbf{14}&\textbf{34}&\PositiveOps\\
Source 58 & $\PP$ &15&24&90\\
Source 59 free-coordinate form & $\PP$ &9&1,910,404&473\\
New degree tradeoff & $\PP$ &9&\NineDegree&\NineOps\\
\bottomrule
\end{tabular}
\caption{Separate complexity measures. The unchanged degree-24 point is
included to prevent a misleading claim that degree 34 is best at every
positive variable count. The new nine-variable point has its own gate cost.}
\label{tab:comparison}
\end{table}

The 42-integer result reduces the number of unknowns by $3/45=6.67\%$
at the same degree. It also reduces the quadratic residual count from 28
to 27 and the number of expanded terms with the four coefficient ports
left symbolic from 407 to 391. The 14-positive result lowers degree by
$18/52=34.62\%$, paying \PositiveExtra\ extra arithmetic operations.
The previously supplied 66-operation, degree-157, 15-positive polynomial
remains a separate useful tradeoff \cite{p51}.

\begin{theorem}[Closed consistency equations]
\label{thm:main}
Bind the four ports of $Q_{42}$ and $P_{14}$ to the complete source-34
tuple in Section~\ref{sec:coefficients}. Then
\begin{align*}
 \ZFC&\vdash
 \bigl[\exists\mathbf x\in\ZZ^{42}\;Q_{42}(\mathbf x)=0\bigr]
 \longleftrightarrow\exists p\in\NN\;\refute(p),\\
 \ZFC&\vdash
 \bigl[\exists\mathbf x\in\PP^{14}\;P_{14}(\mathbf x)=0\bigr]
 \longleftrightarrow\exists p\in\NN\;\refute(p).
\end{align*}
Consequently their respective absence-of-solutions statements are
provably equivalent in ZFC to $\Con(\ZFC)$.
\end{theorem}

Here $\refute$ is the syntactic refutation predicate used by the inherited
compiler's calibrated ZFC endpoint. The theorem is an ordinary mathematical
composition, formalizable in ZFC. It is not a report that a new Lean, Rocq,
or object-language ZFC derivation has been compiled. Section~\ref{sec:boundary}
identifies that distinction precisely.

There is no single ordering of ``simplest'' that minimizes unknowns, degree,
coefficient height, number of monomials, and arithmetic circuit size at once.
For unrestricted integers and low degree, $Q_{42}$ is the main delivered
advance. For a more compact positive-domain display, $P_{14}$ is a useful
balanced choice. The nine-positive polynomial belongs to the smallest-variable
part of the inspected frontier.

\section{The complete 42-integer quartic}
\label{sec:quartic}

\subsection{Variables, abbreviations, and residuals}

The following are the \emph{42 supplied unknowns}, all unrestricted integers:
\begin{align*}
&J,F,\alpha,z,f,h,r,v,w,\tau,\eta,\zeta,y,\gamma,o,j,\\
&N,E,P,c,a,X,D,C_2,L,F_2,V,g,I,B_{\!a},\Delta,u,\rho,\\
&a_1,a_2,a_3,\ b_1,b_2,b_3,\ d_1,d_2,d_3.
\end{align*}
The names $v,\rho,a_i,b_i,d_i$ correspond respectively to
\texttt{v\_scale}, \texttt{q\_half\_root},
\texttt{FA\_pos}$i$, \texttt{FC\_pos}$i$, and
\texttt{EZ\_pos}$i$ in the executable data. The subscript on $B_{\!a}$
distinguishes an auxiliary variable from a fixed coefficient radix.

Define the affine polynomial abbreviations
\[
 q=\beta J+1,\quad k=\eta+\zeta,\quad C=q-F-\alpha,
 \quad H=4a+3,\quad s=2v+1,\quad i=2r+1,
\]
and
\[
 \Pos(x,y,t)=1+x^2+y^2+t(t+1).
\]
Only these six affine quantities are implicit in the following residual
list. In particular, $N,E,P,c,a,X,D,C_2,L,F_2,V,g,I,B_a,\Delta$
are supplied coordinates, constrained by the listed equations, rather than
unpaid intermediate circuit values.

\begin{alignat}{2}
 r_1&=L-q^2+1,&\qquad r_2&=N-q(L+1),\nonumber\\
 r_3&=X-qw,&r_4&=a-E-sN,\nonumber\\
 r_5&=E-X(a-E),&r_6&=c-k(a-E)-\eta,\nonumber\\
 r_7&=P-E(c-\eta)-k,&r_8&=\tau^2-P(P-k)-1,\nonumber\\
 r_9&=D-X-ac-\gamma H,&r_{10}&=\Delta-a^2-H,\nonumber\\
 r_{11}&=C_2-c^2,&r_{12}&=D^2-\Delta C_2-1.
 \label{eq:quartic-first}
\end{alignat}
The next two residuals contain the packing and transport data:
\begin{align}
 r_{13}&=k-1-hE-\bigl[(q-F)(N-q)-CL+(M_C+qM_F)J\bigr],\nonumber\\
 r_{14}&=(K_0+w)C-F-z(q-1).
 \label{eq:quartic-packing}
\end{align}
The remaining arithmetic residuals and the four guards are
\begin{alignat}{2}
 r_{15}&=F_2-f^2,&\qquad r_{16}&=V+c-fo,\nonumber\\
 r_{17}&=V-cj+k-1-hE,&r_{18}&=g+V^2-y^2,\nonumber\\
 r_{19}&=I-iC_2,&r_{20}&=B_a-I^2,\nonumber\\
 r_{21}&=B_a-\Delta(F_2-1),&r_{22}&=y^2-B_ag-1,\nonumber\\
 r_{23}&=w-4u^2,&r_{24}&=q-2\rho^2,\nonumber\\
 r_{25}&=F\alpha-\Pos(a_1,a_2,a_3),&
 r_{26}&=FC-\Pos(b_1,b_2,b_3),\nonumber\\
 r_{27}&=\eta\zeta-\Pos(d_1,d_2,d_3).&&
 \label{eq:quartic-last}
\end{alignat}
The requested single polynomial is exactly
\begin{equation}
 \boxed{\displaystyle Q_{42}=\sum_{\nu=1}^{27}r_\nu^2=0.}
 \label{eq:quartic}
\end{equation}
All 27 residuals are quadratic in the supplied unknowns. The polynomial
is therefore quartic. It has exact degree four because $\tau$ appears
only in $r_8$, as $\tau^2$, and hence
\[
 [\tau^4]Q_{42}=1.
\]
This remains true after the large fixed ports are replaced by their actual
integers. The exact residual and expanded symbolic polynomial data are
\path{quartic/quartic42.json} and
\path{quartic/quartic42_expanded.json}.

\subsection{What changed from source 57}

The first 22 arithmetic residuals are unchanged from the literal source-57
quartic. Its square constraint $w=u^2$ becomes $w=4u^2$. Its two separate
equations $\eta=\Pos(\cdot)$ and $\zeta=\Pos(\cdot)$, with six private
coordinates, are replaced by $r_{27}=0$, with three private coordinates.
The dyadic and simplex guards are unchanged. Thus
\[
 45-6+3=42\quad\hbox{unknowns},\qquad 28-2+1=27\quad\hbox{residuals}.
\]
The restriction to an even square root costs no additional unknown and
does not increase degree. Its justification concerns the canonical
accepting witnesses; it is not a claim that every zero of the old equation
has an even square root at the same coordinates.

The source-57 prose uses an unshifted transport formula in one passage,
while its literal row 14 is the shifted formula \eqref{eq:quartic-packing}.
Our construction preserves the actual polynomial, and our proof uses
\[
 (K_0+w)C=F+z(q-1),\qquad z_{\rm old}=z+1.
\]
The symbolic checker explicitly verifies this identity and rejects the
off-by-one alternative. This correction does not change the source data
or the source's intended existence theorem.

\section{Why one product guard can replace two sign guards}
\label{sec:sign}

\subsection{Elementary guards and packing bounds}

\begin{lemma}[Positive ternary representation]
For all integers $x,y,t$, $\Pos(x,y,t)>0$. Every positive integer is
represented by $\Pos$.
\end{lemma}
\begin{proof}
The integer $t(t+1)$ is nonnegative. For a positive integer $n$, the number
$4n-3$ is $1$ modulo $4$, so it is not of the form $4^a(8b+7)$.
Legendre's three-square theorem represents it as a sum of three squares.
Exactly one square is odd; writing it as $(2t+1)^2$ and the other two
as $(2x)^2,(2y)^2$ gives the required representation.
\end{proof}

At an integer zero of \eqref{eq:quartic}, every residual vanishes.
The authentic fixed coefficient has $\beta\ge15$. From $q=2\rho^2$
and $q=\beta J+1$, $J\le-1$ is impossible by sign, and $J=0$ would
give $1=2\rho^2$. Thus $J>0$ and $q\ge16$.
The two product guards give $F\alpha>0$ and $FC>0$. The integers
$F,\alpha,C$ have the same sign and sum to $q>0$, so all three are
positive.

Restore the abbreviations
\begin{align*}
Y&=sq^3,&X&=wq,&E&=XY,&a&=(X+1)Y,\\
G&=q(q-F)-C,&
R&=G(q^2-1)+(M_C+qM_F)J.
\end{align*}
The first seven rows justify these identities over arbitrary integer
assignments. For example $N=q^3$, $a-E=sN=Y$, and $r_5=0$ gives
$E=XY$. Row 13 is exactly $k=R+1+hE$ because
\[
 (q-F)(N-q)-CL=(q^2-1)[q(q-F)-C].
\]
The exact coefficient bounds are
$0<M_C<\beta$ and $\beta<M_F<2\beta$.
Since $F,\alpha,C>0$ and their sum is $q$,
\[
 2q-1\le G\le q^2-q-1.
\]
Therefore, without using the sign of $z$ or $s$,
\begin{equation}
 q(q^2-1)<R<q^4-q^3<q^4-2,
 \qquad R>3q+1.
 \label{eq:Rbounds}
\end{equation}
For the upper bound, use
$(M_C+qM_F)J<(1+2q)(q-1)$ and expand.

\subsection{The extra modulus margin}

The new guard gives $\eta\zeta>0$, but it has not yet proved that
either coordinate is positive. Write
\[
 \delta=\sgn(\eta)=\sgn(\zeta),\quad
 K=|\eta|+|\zeta|,\quad \eta_0=|\eta|.
\]
Then $\delta\in\{-1,1\}$, $K\ge2$, $0<\eta_0<K$ and $k=\delta K$.
The odd integers $s=2v+1$ and $i=2r+1$ are nonzero.

Initially $w=4u^2\ge0$. If $w=0$, then $X=0$, $a=Y$, and
$D=ac+\gamma H$. Reducing the main norm
$D^2-\Delta c^2=1$ modulo $H$, with $\Delta=a^2+H$, gives
$0\equiv1\pmod H$. But $|Y|\ge q^3$ implies $|4Y+3|>1$.
Thus $w\ne0$, and consequently
\begin{equation}
 w\ge4,\quad X\ge4q,\quad
 M:=|E|=X|Y|\ge4q^4>3R+1.
 \label{eq:modulus}
\end{equation}
This is the decisive strengthening. With only $w=u^2$, the immediate
bound would be $M\ge q^4$, insufficient for the sign argument below.

\subsection{Pell indices before either sign is recovered}

Let $Z=|Y|$, $B=|a+2|$, $P_0=2XY^2+1$, and $c_0=|c|$.
We have $B>M>R$, $\Delta=B^2-1$, and $P_0>B\ge2$.
The two possible signs of $Y$ give respectively
\begin{align*}
 Y>0&:\quad KZ<c_0<K(Z+1),\\
 Y<0&:\quad K(Z-1)<c_0<KZ.
\end{align*}
In particular $c_0>K$ in both cases.

Write $\psi_A(n),\chi_A(n)$ for the integer Pell sequences defined by
\[
 (A+\sqrt{A^2-1})^n=\chi_A(n)+\psi_A(n)\sqrt{A^2-1}.
\]
The first norm in row 8 is
\[
 \tau^2-XY^2(XY^2+1)k^2=1.
\]
The least positive Pell unit for $v(v+1)$ has second coordinate $2$,
so the first and main norms supply indices $n,p\ge1$ with
\[
 K=2\psi_{P_0}(n),\qquad
 c_0=\psi_B(p),\qquad |D|=\chi_B(p).
\]
Since $P_0>B$ and $c_0>K$, monotonicity gives $p>n$.
The recurrence $\psi_{P_0}(n)\equiv n\pmod M$ follows from
$P_0\equiv1\pmod M$. Together with the retained index row it gives
\begin{equation}
 2\delta n\equiv R+1\pmod M.
 \label{eq:sign-congruence}
\end{equation}
If $\delta=1$, then $n\ge(R+1)/2$.
If $\delta=-1$, then
\begin{equation}
 n\ge\frac{M-R-1}{2}>R.
 \label{eq:negative-rank}
\end{equation}
In both branches $n\ge24$ and $p\ge25$. Elementary Pell growth now gives
\begin{equation}
 c_0>B\Delta^2,\qquad c_0>2p,\qquad c_0>2R.
 \label{eq:rank-bounds}
\end{equation}
For example $\psi_B(p)\ge(2B-1)^{p-1}$ proves the first two
inequalities; $c_0\ge\psi_B(2)=2B>2R$ proves the third.

\subsection{The auxiliary equation fixes the main index}

The unchanged quadratic rows assert, with $\rho_0=|i|c_0^2$,
\begin{gather*}
 \rho_0^2=\Delta(f^2-1),\qquad
 \rho_0^2V^2-(\rho_0^2-1)y^2=1,\\
 V=of-c=jc-R.
\end{gather*}
The ordinary relaxed-rank theorem, applied to the positive pair $(B,c_0)$
using \eqref{eq:rank-bounds}, gives
\[
 |f|=\chi_B(m),\qquad p\mid m,\quad c_0\mid m,\quad m>2p.
\]
Its proof is retained in \path{upstream/PELL_RELAXED_AUXILIARY_PROOF.md}.
It uses squares of $c$ and absolute values, so it does not assume the sign
that we are trying to prove.

For clarity, the signed-index step can be traced explicitly. The auxiliary
Pell classification supplies an odd index $\ell=2j_0+1$ and a sign $\sigma$
with
\[
 V=\sigma\mathcal Q_{j_0}(\rho_0^2),
\]
where the integral odd Chebyshev quotient satisfies
\[
 \mathcal Q_j(1-B^2)=(-1)^j\psi_B(2j+1),\qquad
 \mathcal Q_j(0)=(-1)^j(2j+1).
\]
The exposed quotient equations give
$V\equiv-c\pmod{|f|}$ and $V\equiv-R\pmod{c_0}$.
As $\rho_0^2\equiv1-B^2\pmod{|f|}$, squaring the first congruence and
using Pell duplication yields
\[
 \chi_B(2\ell)\equiv\chi_B(2p)\pmod{\chi_B(m)}.
\]
The signed step-down lemma, with $0<2p<m$, gives
$\ell\equiv p$ or $-p\pmod m$. The second congruence and $c_0\mid m$
therefore imply $R\equiv p$ or $-p\pmod{c_0}$.
The strict bounds in \eqref{eq:rank-bounds} force
\begin{equation}
 p=R.
 \label{eq:pR}
\end{equation}
The quotient and step-down identities are the same exact inherited lemmas
used by the signed source; their statements and proofs are included in
\path{upstream/SIGNED_DOMAIN.md} and
\path{upstream/EXPLORATION_FIXED_MINUS_INDEX_PARITY.md}.

If $\delta=-1$, equations \eqref{eq:negative-rank} and \eqref{eq:pR}
contradict $p>n$. Hence $\delta=1$, proving
\[
 \eta>0,\qquad\zeta>0.
\]
This completes the new sign-recovery argument. Each positivity hypothesis
has been obtained before the theorem requiring it is invoked.

\subsection{Restore the old zero and prove the converse}

Given a new zero, replace its square witness $u$ by $u_{\rm old}=2u$.
Supply ternary representations of the now-positive $\eta$ and $\zeta$.
All other coordinates and all four coefficient ports are unchanged.
This is a zero of the literal source-57 quartic. Its inherited signed
compiler theorem decodes a refutation.

For completeness, start with a refutation and the source compiler's
canonical accepting witness. Its radix is $q=2^t$ with $t$ odd, and its
rank has $R\equiv3\pmod4$. Moreover
\[
 w=2^{R-t},\qquad R>t+2.
\]
Thus
\[
 u=2^{(R-t-2)/2}
\]
is an integer and gives $w=4u^2$. Retain the dyadic and simplex witnesses,
and choose a ternary representation of the positive number $\eta\zeta$.
Every new residual vanishes. This proves both directions of the new
existence equivalence.

For the imported source's transport conclusion, the literal shifted
formula gives $z>0$ after its power recovery: $X=2^R$ implies
$w=X/q>q>F$, whence $(K_0+w)C-F>0$. Therefore
$z(q-1)>0$. This verifies the needed domain assertion without the
off-by-one formula mentioned above.

\section{A degree-34 equation in fourteen positive unknowns}
\label{sec:positive14}

A second construction improves a different part of the measured frontier.
It has fourteen strictly positive unknowns, exact degree $34$, and $87$
live binary arithmetic operations. The inspected fourteen-positive
degree-$52$, $79$-operation forms remain lower-cost alternatives. Thus the
new point reduces degree by $34.62\%$, at a cost of eight extra operations;
it does not reduce the nine-positive universal-variable bound.

This construction uses the same whole source-34 coefficient tuple as the
main result: $\beta=2^{1064921}-1$, together with its authentic $K_0,M_C,M_F$.
The high mask $M_F$ is already shifted. The first-root carrier from source
58 and the index modulus from source 57 are inherited ingredients. The new
combination uses a single positive main-root carrier to eliminate the index
quotient $h$. Keeping the normalized strong equation avoids the
large-ordinate hypothesis needed by the ordinary-strong variant.

\subsection{The explicit polynomial}

Supply precisely
\[
 J,F,\alpha,z,f,i,T,w,Y,r,\eta,\zeta,u,t\in\mathbb Z_{>0}.
\]
All subsequent names are polynomial abbreviations, not extra unknowns:
\begin{align*}
 q&=\beta J+1,& C&=q-F-\alpha,&X&=wq,&E&=XY,\\
 k&=\eta+\zeta,&c&=kY+\eta,&a&=E+Y,&H&=4a+3,\\
 \Delta&=a^2+H,&b_0&=X+ac,&\ell&=EkY,&I&=ic^2.
\end{align*}
Define
\begin{align*}
 R&=[(q-1)(q-F)+\alpha](q^2-1)+(M_C+qM_F)J,\\
 M&=4q^4,\qquad e=k-R-1,\qquad Z=Mt-e,\\
 W&=Mb_0+HZ,\qquad Q=\Delta(f^2-1),\qquad
 V=c(Tf-1)-Rf^2.
\end{align*}
The five residuals are
\begin{align}
 A_0&=[q^3(\ell+r)+Y]^2-q^6[\ell(\ell+k)+1],\label{eq:p14-first}\\
 A_1&=W^2-M^2(1+\Delta c^2),\label{eq:p14-main}\\
 A_t&=(K_0+w)C-F-z(q-1),\label{eq:p14-transport}\\
 A_s&=(f^2-1)-\Delta I^2,\label{eq:p14-strong}\\
 A_u&=(V+u)^2-Qu(2V+u)-1.\label{eq:p14-aux}
\end{align}
Our second single equation is
\begin{equation}
 \boxed{P_{14}=A_0^2+A_1^2+A_t^2+A_s^2+A_u^2=0.}
 \label{eq:p14}
\end{equation}
Only fixed powers and integer coefficients occur. The circuit is
\path{positive14/pell14_degree34.json}. Its 78-operation residual evaluator
becomes an 87-operation single-polynomial evaluator after five squarings
and four additions.

\subsection{Soundness and the order of index recovery}

At a positive zero, all five residuals vanish. Equation
\eqref{eq:p14-transport} gives $C>0$, because $K_0+w>0$ and
$F+z(q-1)>0$. The old transport coordinate is $z_{\rm old}=z+1>0$.
Substituting it in the old residual
$(K_0+w)C+q-F-z_{\rm old}(q-1)-1$ gives exactly $A_t$,
recovering the authentic transport equation.
The inherited outer bounds consequently give
\[
 q\ge16,\qquad 3q<R<q^4-2.
\]

Equation \eqref{eq:p14-first} says that
$\ell+r+Y/q^3$ has an integral square. A rational square root of an integer
is integral; hence $q^3\mid Y$. Set
\[
 s=Y/q^3>0,\qquad \tau=\ell+r+s>0.
\]
Then $\tau^2=\ell(\ell+k)+1$. Write $P_0=2XY^2+1$ and $A=a+2$.
First-Pell classification gives
\[
 \tau=\chi_{P_0}(n),\qquad k=2\psi_{P_0}(n),\qquad n\ge1.
\]
The case $n=1$ would give $\tau=\ell+1$, whereas the retained gap gives
$\tau\ge\ell+2$. Thus $n\ge2$, and
\begin{equation}
 k\ge4P_0>8XY^2,\qquad
 c>kY>8XY^3\ge8q^{10}>2R,\qquad c>X.
 \label{eq:p14-earlybounds}
\end{equation}
This establishes the needed large ordinate \emph{before} restoring $h$.

The main carrier is positive. Indeed $a\ge q^4\ge3$ gives $H\le5a$,
so $b_0/H\ge c/5>k$. Consequently
\[
 W/H=Mb_0/H+Mt-k+R+1>0.
\]
By \eqref{eq:p14-main}, $D=W/M$ is an integral rational square root,
and therefore a positive integer satisfying $D^2-\Delta c^2=1$.
In particular,
\begin{equation}
 M\mid He.\label{eq:p14-divisibility}
\end{equation}
Neither $e/M$ nor $e/E$ has yet been assumed integral.

Main-Pell classification gives $D=\chi_A(p)$ and $c=\psi_A(p)$.
Since $P_0>A$ and $c>k$, monotonicity gives $n<p$; hence $p\ge3$
and $c>2p$. The normalized strong equation gives
\[
 f=\chi_A(m),\qquad I=\psi_A(m)=ic^2.
\]
Strong divisibility gives $p\mid m$. If $m=pj$, the composition formula
modulo $c$ reads
\[
 \frac{\psi_A(pj)}c\equiv j\chi_A(p)^{j-1}\pmod c.
\]
The left side is divisible by $c$ and $\chi_A(p)$ is invertible modulo $c$.
Thus $c\mid j$, so $pc\mid m$, $c\mid m$, and $m>2p$.

Put $\rho=\Delta I>1$. Equation \eqref{eq:p14-strong} implies
$Q=\rho^2$, and \eqref{eq:p14-aux} becomes
\[
 \rho^2V^2-(\rho^2-1)(V+u)^2=1.
\]
Neither $V$ nor $V+u$ vanishes. The literal formula for $V$, together with
$f^2\equiv1\pmod c$, gives
\[
 V\equiv-c\pmod f,\qquad V\equiv-R\pmod c.
\]
The inherited signed auxiliary-index lemma now applies with $0<2p<m$.
For precision, its odd-index quotient identities are
\[
 Q_j(1-A^2)=(-1)^j\psi_A(2j+1),\qquad
 Q_j(0)=(-1)^j(2j+1).
\]
They yield an auxiliary index congruent to $\pm p$ modulo $m$ and to
$\pm R$ modulo $c$. Since $c\mid m$ and $0<p,R<c/2$, one obtains
\begin{equation}
 p=R.\label{eq:p14-rank}
\end{equation}
No assertion that the deleted quotient exists has entered this argument.

We can now use \eqref{eq:p14-divisibility}. As $Y=sq^3$, the odd integer
$H$ is congruent to $3$ modulo $q$. If $3\mid q$ then $27\mid Y$ and
$H\equiv3\pmod{27}$, so
\[
 \gcd(M,H)\le3,\qquad
 \mu=M/\gcd(M,H)\ge4q^4/3>R.
\]
The first recurrence gives $k\equiv2n\pmod{4XY^2}$, and $4XY^2$ is
divisible by $M$. Therefore $\mu\mid2n-R-1$. Since $1\le n<p=R$,
$|2n-R-1|<R<\mu$, and hence $2n=R+1$. It follows that
$e=k-2n$ is a positive multiple of $4XY^2$. Thus
\[
 h=e/E>0,\qquad \gamma=t-e/M\in\mathbb Z
\]
are restored, with $k-hE-R=1$ and $D=b_0+\gamma H$.
Finally $D>(a+1)c$ and \eqref{eq:p14-earlybounds} give
$D-b_0>c-X>0$, so $\gamma>0$.

Taking $y_{\rm old}=|V+u|>0$ restores the complete positive normalized
source system. Its inherited compiler and syntactic calibration therefore
give a finite ZFC refutation. The separate mathematical review also notes
that $V+u$ is already positive: if it were nonpositive, then
$V=V+u-u<V+u\le0$ would make the auxiliary left side at least $Q>1$.
The absolute value is harmless and avoids depending on that strengthening.

\subsection{Completeness: why both new coordinates are positive}

Start from a canonical positive normalized compiler witness for a finite
refutation. It has $p=R=2n-1$, $R\equiv3\pmod4$, $n\ge3$, and $k>XY$.
Its first recurrence makes $e=k-R-1>0$ divisible by $M$. Define
\[
 t=\gamma+e/M>0.
\]
This gives $W=MD$. The plus sign in the definition of $t$ is essential;
subtracting $e/M$ from the old positive $\gamma$ would not be justified.

For the first root define $r=\tau-\ell-s$. Since
$\tau^2=\ell^2+\ell k+1$, $\ell>k$, and $\tau<\ell+k$,
\[
 \tau-\ell=\frac{\ell k+1}{\tau+\ell}>\frac{k}{3}
 >\frac{XY}{3}>\frac{Y}{q^3}=s.
\]
Hence $r>0$ and \eqref{eq:p14-first} holds.

Retain the canonical normalized $f,i$ and set $\rho=\Delta I$.
Choose fresh auxiliary witnesses
\[
 V_* =\chi_\rho(R)/\rho,\qquad y_* =\psi_\rho(R).
\]
Odd $R$ makes $V_*$ integral. The odd-index congruences, with
$R\equiv3\pmod4$, give $V_*\equiv-c\pmod f$ and
$V_*\equiv-R\pmod c$. As $\gcd(c,f)=1$ and $f^2\equiv1\pmod c$,
\[
 T=\frac{V_*+c+Rf^2}{cf}>0
\]
is an integer. The Pell identity and $V_*>1$ imply
\[
 y_*^2-V_*^2=\frac{V_*^2-1}{\rho^2-1}>0.
\]
Thus $u=y_*-V_*>0$ is integral and \eqref{eq:p14-aux} holds.

The final coordinate change is the established positive transport shift.
Canonical $X=2^R$ and $R>3q$ give $w=X/q>q>F$. The source transport
equation is
\[
 (K_0+w)C=F+(z_{\rm old}-1)(q-1).
\]
Its left side exceeds $F$, so $z=z_{\rm old}-1>0$.
All fourteen supplied variables are positive and all residuals vanish.
This completes both existence directions. Negation and the inherited
calibration prove in ZFC the required consistency equivalence for $P_{14}$.

\subsection{Degree, implementation, and limits}

Two exact identities expose the cancellations:
\begin{align*}
 A_0={}&q^6[r^2+\ell(2r-k)-1]+2q^3(\ell+r)Y+Y^2,\\
 A_1={}&M^2[X^2+2acX-Hc^2-1]+2Mb_0HZ+H^2Z^2.
\end{align*}
They give the respective residual degrees $12,17,2,16,15$.
Only $A_1^2$ reaches degree $34$. Its leading homogeneous part is
\[
 \operatorname{in}_{34}(P_{14})
 =2^{14}(\beta J)^{20}w^4Y^6(\eta+\zeta)^2t^2.
\]
In particular the coefficient of $J^{20}w^4Y^6\eta^2t^2$ is
$2^{14}\beta^{20}>0$, proving exact degree.

The literal primary evaluator has 48 multiplications, 23 additions and 16
subtractions. The input-plus-one translation has fourteen nonnegative
unknowns, the same degree, and 101 operations. A transparent unshifted
transport comparator has 88 operations. The separate cancellation graph
is a degree certificate, not the evaluator used for the 87-operation claim.

The finite verifier checks all four actual coefficient payloads against
their pinned hashes, exact local polynomial identities, liveness and counts,
the universal degree upper bound, degree-attaining finite-field polynomial
rays, and independently evaluated literal formulas. These finite checks
support the implementation. The unbounded equivalence is the mathematical
argument above with its inherited Pell, compiler and calibration inputs;
no new end-to-end kernel certificate is asserted.


\section{A deflated majorant in nine positive variables}
\label{sec:deflated-nine}

This contribution keeps the source and free-coordinate compiler of
\path{59-degree-separated/nine_free_adjunct} at ProveIt commit
\texttt{beb0192cb8dc36f78c369c080f2adfbc1887a486}.
It changes the polynomial that strictly majorizes a radical in the last
norm equation.  The resulting literal polynomial has nine strictly positive
unknowns, exact total degree \(1{,}573{,}266\), and a division-free evaluator
with \(504\) binary operations.  The predecessor has degree
\(1{,}910{,}404\) and \(473\) operations: the degree decreases by
\(337{,}138\), or approximately \(17.6475\%\), at the cost of \(31\)
additional operations.  Thus this is a degree/size tradeoff.

\subsection{New algebraic ingredient}

Use the inherited registers
\begin{align*}
 J&=4A-5,&D&=(A^2-1)C^2+1,&e&=iJC,\\
 E&=eDC,&F&=(A^2-1)E^2+1,&L_m&=CD+1-A,\\
 G&=A+FL_m,&H&=C+F(2R+1+(j-1)C),&I&=(G^2-1)H^2+1.
\end{align*}
The required domain is \(A\geq3\), \(C,i\geq1\).  Set
\begin{align*}
 N_*&=2G+L_m,&N_1&=2e^2D(D-1)L_m+3C,\\
 N_2&=2e^2(D-1)L_m,&N_3&=2e^2(L_m-C),&N_4&=2e^2C.
\end{align*}
The new polynomial is
\begin{align}
 T_7={}&e\bigl[(2048D^2-1024D-256)N_*
       -128N_1-80N_2-56N_3-42N_4 \notag\\
       &\hspace{22mm}+i^2(1568-448A+288C)\bigr]+640i+4096.
       \label{eq:deflated-T7}
\end{align}
No displayed register adds a quantified unknown.  They are abbreviations
for ordinary integer polynomials in the nine final inputs.

The simplifying relation is
\[
 e^2=16i^2D-i^2\bigl((40A-41)C^2+16\bigr).
\]
It cancels a complete high-degree part of the quotient remainder used by
the earlier majorant.  To make this exact, define
\begin{align*}
 U_6={}&e\bigl[(1024D^2-512D-128)N_*
             -64N_1-40N_2-28N_3-21N_4\bigr],\\
 T_6={}&U_6-224i^2e(A-1-3C)+560i^2e+320i+2048.
\end{align*}
Then \(T_7=2T_6-1056i^2eC\).  The variable-dependent subtraction is
justified by the following exact certificate.

Put
\begin{align*}
 t_7={}&2048D^7-1024D^6-256D^5-128D^4
           -80D^3-56D^2-42D-33,\\
 d_7={}&109824D^6+36608D^5+18304D^4+9984D^3
           +5460D^2+2772D+1089,\\
 R_7={}&D^5T_7-et_7N_*.
\end{align*}
Exact expansion gives
\[
 t_7^2-2048^2D^{13}(D-1)=d_7>0.
\]
Since the sum of the absolute negative coefficients of \(t_7\) is
\(1619<2048\), one has \(t_7>0\) for \(D\geq1\).  The delivered
file \texttt{data/order7\_majorant\_positive\_certificate.json} gives
the complete exact identity
\begin{equation}
 R_7-4096D^5
   =\sum_{(u,v,w)\in\mathcal S} a_{uvw}(A-3)^u(C-1)^v i^w,
 \qquad |\mathcal S|=264,
 \label{eq:deflated-positive-certificate}
\end{equation}
where every \(a_{uvw}\) is a strictly positive integer.  At \(A=3,C=1\)
the right side is
\(49{,}748{,}823i+293{,}313{,}636i^3>0\).
The checker reconstructs every coefficient from
\eqref{eq:deflated-T7}; this is a finite polynomial identity and a uniform
positivity proof, rather than a test on sampled parameter values.

For completeness, write \(X=e^2D^3(D-1)\).  Then
\(DF=X+D\) and \(G=(X/D)L_m+CD+1\).  The elementary inequality
\(\sqrt{X+D}<\sqrt X+D/(2\sqrt X)\) gives
\[
 4096\sqrt{DF}\,G
 <2048eD^2\sqrt{1-1/D}\,N_*+4096v,
 \qquad
 v=\frac{CD+1}{2eD\sqrt{1-1/D}}<1.
\]
The final inequality follows from \(D\geq A^2\geq9\) and \(e\geq7C\).
Combining it with \eqref{eq:deflated-positive-certificate} and the square
identity proves the unconditional bound
\begin{equation}
             T_7>4096\sqrt{DF}\,G>0.
             \label{eq:deflated-majorant}
\end{equation}
The divisions occur only in this proof, not in the polynomial.

\subsection{Final polynomial and semantic equivalence}

Keep the inherited positive source parameters and normalized margin:
\begin{align*}
 B&=2^{2134939}b^2,&N&=2^{17079501}bB^{1201},\\
 t&=bB^{116},&m&=K^2(4t-g)-16t(C-KY)^2,\\
 \Omega&=Jm,&E_0&=J\bigl(n+(bw)^2\bigr)+Q.
\end{align*}
Set
\[
\begin{aligned}
 a_0&=(P^2-1)K^2+1,& a_1&=4096^2DFI,\\
 V_0&=PK,& V_1&=T_7H,\\
 \Theta&=E_0-\Omega(V_0+V_1).
\end{aligned}
\]
The equation is
\begin{equation}
 F_9=\bigl(\Theta^2-\Omega^2(a_0+a_1)\bigr)^2
            -4\Omega^4a_0a_1=0.
 \label{eq:deflated-final}
\end{equation}
Its only unknowns are \(b,g,h,s,w,C,i,j,n\), all strictly positive.
The source coefficients, source positions and ratio conditions are exactly
those of the pinned predecessor.  The four full native coefficient
payloads are supplied; there is no unspecified program-index parameter.

Equation \eqref{eq:deflated-majorant} gives \(V_1>\sqrt{a_1}\), and
\(V_0^2-a_0=K^2-1>0\).  Also
\(E_0=Jn+(J-2)(bw)^2+3bwC+2>0\).  At a zero, \(m=0\) is impossible.
Real factorization yields roots \(r_0^2=a_0,r_1^2=a_1\) with
\[
 n+(bw)^2+Q/J=m(V_0+V_1+r_0+r_1).
\]
Both the left side and the right-hand parenthesis are strictly positive.
Consequently \(m>0\) follows before any integrality assertion.  The
inherited margin argument recovers \(C>2R+1+bw\) and \(C>K\), hence
\(a_1>a_0\).  The rational sum \(r_0+r_1\) is therefore nonzero, so
\(2(r_0+r_1)r_0=(r_0+r_1)^2+a_0-a_1\) makes the roots rational and
then integral.  The same rational-square argument makes \(r_1/4096\)
integral, and the factor equation gives \(J\mid Q\).  These are the
unchanged ratio hypotheses.

Conversely, take the inherited ratio witnesses and positive roots, and put
\(n=m(V_0+V_1+r_0+r_1)-(bw)^2-Q/J\).  It is an integer, and it is
positive because \((bw)^2+Q/J<2D\), \(m\geq1\), and
\(V_1>4096D\).  It makes \eqref{eq:deflated-final} vanish.  Thus the
positive zero set is equivalent to the pinned source zero set.  Composing
with its native compiler and K215 calibration gives, in ZFC, absence of
positive zeros if and only if \(\operatorname{Con}(\mathrm{ZFC})\).
This is an ordinary mathematical proof with named inherited dependencies;
the new semantic equivalence is not claimed to have been kernel checked.

\subsection{Exact degree and finite verification}

Let \(\delta_j=V_j^2-a_j\) and
\(B_*=E_0^2-2E_0\Omega(V_0+V_1)+\Omega^2(\delta_0+\delta_1)\).
The two-defect identity gives
\[
 F_9=B_*^2+4\Omega^2V_0V_1B_*
          +4\Omega^4(\delta_0V_1^2+a_0\delta_1).
\]
The exact majorant defect identity is
\begin{align*}
 D^{10}(T_7^2-4096^2DFG^2)
  ={}&e^2d_7N_*^2
     +2048^2D^{10}\bigl[XL_m^2-4DG(CD+1)\bigr]\\
    &+2et_7N_*R_7+R_7^2.
\end{align*}
Here \(\deg R_7=187305\), with highest weighted term
\(2880A^{10}C^{12}i^3\).  The mixed term uniquely dominates the last
identity.  It yields
\[
 \deg(T_7^2-4096^2DFG^2)=280959,\qquad
 \deg\delta_1=598873,\qquad \deg B_*=786633.
\]
The square \(B_*^2\) uniquely dominates the final two-defect expansion,
proving \(\deg F_9=1{,}573{,}266\).  A separate circuit for
\(D^{20}F_9\) supplies an exact nonzero diagonal leading coefficient
and the matching total-degree upper bound.  Its extra operations are not
included in the primary evaluator count.

Running \texttt{python3 verify.py} checks all four native byte payloads,
the source isolation data, the 264-term positivity identity, the defect
identity, every delivered circuit against its deterministic builder,
graph liveness, exact degrees, costs and domain translation.  It also
performs 480 comparisons with a separately written literal evaluator.
The check uses only Python's standard library.  These finite checks support
the implementation and identities; the unbounded semantic endpoint uses
the explicitly credited source and compiler theorems.


\section{Exact coefficient data and semantic dependencies}
\label{sec:coefficients}

The primary quartic and 14-positive construction both use the entire
source-34 tuple, rather than exchanging individual ports between
different compiler instances. Let $B=2^{1064921}$. Then
\[
 \beta=B-1,\qquad K_0=D_C+B D_R,\qquad M_C=M_{C,\rm native},
 \qquad M_F=M_{F,\rm native}+B-1.
\]
The archive contains the actual unsigned little-endian bytes for these
four integers and a finite builder from the literal 430 machine actions
and 943 Wang tiles. The high mask is already shifted; no second shift is
applied.

\begin{table}[htbp]\centering\small
\begin{tabular}{@{}lrrr@{}}\toprule
Port & Bits & Raw bytes & Compressed bytes\\\midrule
$\beta$ &1,064,921&133,116&164\\
$K_0$ &1,065,043&133,131&9,771\\
$M_C$ &1,064,921&133,116&270\\
$M_F$ &1,064,922&133,116&261\\\midrule
Total &&532,479&10,466\\\bottomrule
\end{tabular}
\caption{The complete source-34 binding. These fixed integers are leaves,
not additional quantified unknowns. Compressed byte counts describe the
supplied deterministic gzip files, not coefficient heights.}
\end{table}

The builder and independently written decoder have been rerun on the
actual fixed values. They check the transition/tile correspondence,
69,230 finite basis-convolution conditions, all 8,801 coefficient digits,
mask populations, strict carry margins, and the finite arithmetic facts
used by rank tuning. The universal compiler correctness is an inherited
mathematical theorem, not a consequence of those finite checks alone.
The source's coefficient proof and original notices are retained unchanged
under \path{coefficients/} and \path{upstream/licenses/}.

The additional nine-positive component has its own complete binding,
explicitly identified in that component's data. It is never assembled by
mixing one port from the source-34 tuple with another source's transport
or mask port.

\begin{table}[htbp]\centering\small
\begin{tabular}{@{}lccc@{}}\toprule
Dependency & $Q_{42}$ & $P_{14}$ & Nine-positive form\\\midrule
Literal native compiler / calibrated endpoint &yes&yes&yes\\
Ordinary Pell and signed-index lemmas &yes&yes&yes\\
Jones ratio representation &no&no&yes\\
Sun's ten-integer reduction &no&no&no\\
New end-to-end proof-assistant certificate &no&no&no\\\bottomrule
\end{tabular}
\caption{Scope of the delivered proofs. The main low-degree constructions
do not use Jones's universal expression.}
\end{table}

The user's permission to change the computational starting point was
taken seriously. A separate direct proof-calculus route already exists
in the pinned repository; its native enumeration-step graph uses 1,742
witnesses before quartic flattening
\cite{nativegraph}. A different formally verified route through Minsky
machines and FRACTRAN is available in the Coq literature \cite{coqh10}.
Those are valid alternatives, but no smaller explicit closed equation
with a measured complete coefficient binding was obtained from them in
this iteration. The delivered results therefore retain the native
calibration while changing its arithmetic representation.

\section{Reproduction, finite certificates, and proof status}
\label{sec:boundary}

\subsection{What the programs verify}

From the extracted archive root, run
\begin{verbatim}
python3 verify.py
\end{verbatim}
The verification driver checks the inherited source identities, actual
coefficient payloads, deterministic regeneration of the new equations,
literal evaluator counts, symbolic cancellation and exact-degree
certificates, and independently evaluated formulas modulo large primes.
The quartic verifier also runs with Python's assertion optimization enabled:
its checks use explicit exceptions, not removable assertions.

For the quartic, the sparse symbolic expansion proves that the first 22
arithmetic rows are exactly the source rows, the scale change is exactly
$-3u^2$ added to the old scale residual, and the dyadic/simplex rows are
unchanged. All 42 variables are live. The exact expansion has 391 terms
when the four fixed ports remain indeterminates. The coefficient of
$\tau^4$ is exactly one. Independent AST evaluation and sparse expansion
agree in 180 modular comparisons using the actual bound ports.

The finite sign and representation checks are supplementary diagnostics.
For example, they test 392 signed interval cases and 1,590 negative-index
congruence cases. The universal sign theorem is the proof in
Section~\ref{sec:sign}, not an extrapolation from those samples.
The positive and nine-variable component certificates state their own
precise scopes and counts.

Rebuild this report with
\begin{verbatim}
cd report
pdflatex -interaction=nonstopmode -halt-on-error report.tex
pdflatex -interaction=nonstopmode -halt-on-error report.tex
\end{verbatim}
The PDF is a compiled version of the supplied LaTeX source. Recompiling
may change PDF metadata bytes; the archive manifest authenticates the
delivered version.

\subsection{Why the composition is a theorem of ZFC}

The asserted internal provability is not inferred merely from the truth
of an external Lean theorem or from agreement on standard numerals.
Work in ZFC with natural-number codes for formulas, finite derivations,
and finite machine runs. The register compiler and binary simulation
are uniform inductions on such codes and runs. The inherited ordinary
K215 proof explicitly states this calibration in ZFC
\cite{k215}; its exact text is retained as
\path{upstream/K215_CALIBRATION_PROOF.md}.

The calibration's model-theoretic part uses set-sized structures:
Henkin completeness, satisfaction for set structures, a quotient
interpretation with an internally defined well-founded part, and the
constructible submodel of a set model. These are constructions available
in ordinary ZFC. In particular, the well-founded-part interpretation
quantifies over represented subsets, not arbitrary external subsets.
It does not assume a transitive set model, an inaccessible cardinal, or
external well-foundedness. The constructible domain inside a set model
is an ambient subset of that model's domain; no truth predicate for the
ambient universe is required. The exact inspected source interfaces are
listed in \path{reviews/ZFC_INTERNALIZATION_SCOPE.md}.

The arithmetic compiler and the new Pell and polynomial arguments are
likewise uniform coded arguments over the natural numbers. Their
composition can therefore be carried out within ZFC, with quantifiers
ranging over the ambient theory's entire $\omega$. Equivalently, the
same proof runs inside any model of ZFC using that model's naturals;
there is no step converting a nonstandard internal certificate into an
externally finite object. The proof predicate is that of the calibrated
finite ZFC calculus, or a standard presentation connected to it by the
usual primitive recursive proof translations. Negating its existence
gives the stated syntactic consistency sentence.

This ordinary internalization argument explains the meaning of
Theorem~\ref{thm:main}. Exporting it as an actual derivation in the
repository's selected object-language calculus remains a separate
formalization task, exactly as warned in the preserved status document.

\subsection{What is proved and what remains to formalize}

The new work establishes ordinary mathematical existence equivalences.
These use exact inherited theorems about the native compiler, the
calibrated endpoint, and Pell representations. Those mathematical
arguments can be formalized in ZFC, giving the internal consistency
biconditionals in Theorem~\ref{thm:main}. No assumption that ZFC is
consistent enters either direction.

Three different artifacts should be distinguished:
\begin{enumerate}
\item The explicit polynomial and its finite arithmetic certificates,
which are supplied here.
\item An external proof-assistant theorem about the exact new literal
polynomial. No newly compiled end-to-end theorem of this kind is supplied.
\item A formal derivation inside the repository's object-language ZFC
calculus of the desired biconditional. That exported derivation is also
not supplied.
\end{enumerate}
The pinned repository has substantial checked components and separate
closed formal routes. They do not automatically certify a changed
literal equation. The preserved \path{upstream/OBJECT_ZFC_STATUS.md}
records the relevant boundary at the comparison snapshot.

The natural next formalization target for $Q_{42}$ is especially small:
formalize the common-sign bootstrap with the abstract signed-rank lemma,
then prove the exact 27-row adapter. The literal coefficient data and
native compilation theorem can be reused. For $P_{14}$, isolate the
first-root gap's exclusion of index one, the normalized-rank argument,
and the main-root quotient reconstruction before composing them.

The independent reviews bundled with this work are separate AI-assisted
reviews from this research session. They are not represented as external
human peer review or kernel certificates. The archive does not prove
consistency, inconsistency, or any global minimality assertion.

\section{Contribution and remaining opportunities}

The main new quartic idea is reusable beyond this instance. When two
coordinates are known only to have the same sign, a modulus large enough
to separate the negative residue branch can sometimes recover their
individual signs. A cheap restriction on canonical witnesses can provide
that modulus while preserving every accepting computation. Here this
removes three unrestricted integer variables without raising degree.

The positive construction uses a different resource: an existing square
root can carry a rational offset. Its integrality restores a removed
divisibility witness, while its strict positive gap supplies an early
index bound. Both directions matter. Merely eliminating the quotient
from the displayed formula would introduce a circular rank argument;
the proof here obtains the necessary bound first.

Further progress should target an actual measured bottleneck: the
remaining simplex guards in the integer quartic, the degree-17 main
carrier in the positive form, or the source coding and radical-majorant
degrees in the nine-positive form. A new proof calculus or machine may
also reduce the coefficient description, but a smaller state count alone
does not establish a smaller Diophantine equation.

\begin{thebibliography}{99}
\bibitem{collection} ProveIt contributors.
\emph{Con(ZFC) in Diophantine form}, sources 31--59, snapshot \pin.
\url{https://github.com/VladimirReshetnikov/ProveIt/tree/beb0192cb8dc36f78c369c080f2adfbc1887a486/Computability/DiophantineConsistency/con-zfc-diophantine-forms}.
\bibitem{p34} ProveIt research source 34.
\emph{Fourteen unrestricted integers, and a lower-degree fifteen-integer form};
coefficient adapter and quartic50 data. Directory \path{34-integer14/}
in \cite{collection}. Exact copied source identities are in
\path{provenance/SOURCES.json}.
\bibitem{p51} ProveIt research source 51.
\emph{Sparse arithmetic}, including the 66-operation, degree-157 positive
form. Directory \path{51-sparse-arith/} in \cite{collection}.
\bibitem{p55} ProveIt research source 55.
\emph{Congruence separation and positive-word fusion}.
Directory \path{55-congruence-norms/} in \cite{collection}.
\bibitem{p57} ProveIt research source 57.
\emph{Separated degrees and Pell quotients}.
Directory \path{57-separated-degrees/} in \cite{collection}.
\bibitem{p58} ProveIt research source 58.
\emph{Integral-root carriers for Con(ZFC)}.
Directory \path{58-integral-root/} in \cite{collection}.
\bibitem{p59} ProveIt research source 59.
\emph{Degree-separated encodings and an ordinary-root merger for Con(ZFC)}.
Directory \path{59-degree-separated/} in \cite{collection}.
\bibitem{signed} ProveIt research source 36, retained earlier proofs.
\emph{Signed-domain reduction} and \emph{Odd half-binomial scale}.
\path{36-small-tuple/coefficients/source_proofs/} in \cite{collection}.
\bibitem{rank} ProveIt research notes.
\emph{Recovering integrality after relaxing an auxiliary Pell norm};
\emph{Fixed-minus index parity}; \emph{Signed Pell proof}.
\path{Computability/HilbertTenthProblem/Papers/1980/}, snapshot \pin.
\bibitem{k215} ProveIt contributors.
\emph{Ordinary proof of the combined215 machine}, snapshot \pin.
\path{Computability/BusyBeaver/Research/ZFC/candidate-k215/proof.md}.
The verbatim text and exact source hash are included in this archive.
\url{https://github.com/VladimirReshetnikov/ProveIt/blob/beb0192cb8dc36f78c369c080f2adfbc1887a486/Computability/BusyBeaver/Research/ZFC/candidate-k215/proof.md}.
\bibitem{nativegraph} ProveIt contributors.
\emph{Polynomial formula blocks in the native refutation graph},
\path{Computability/DiophantineConsistency/NATIVE_FORMULAS.md}, snapshot \pin.
\url{https://github.com/VladimirReshetnikov/ProveIt/blob/beb0192cb8dc36f78c369c080f2adfbc1887a486/Computability/DiophantineConsistency/NATIVE_FORMULAS.md}.
\bibitem{coqh10} D. Larchey-Wendling and Y. Forster.
\emph{Hilbert's Tenth Problem in Coq (Extended Version)}.
Logical Methods in Computer Science 18(1):35, 2022.
\url{https://arxiv.org/abs/2003.04604}.
\end{thebibliography}
\end{document}
